% ═══════════════════════════════════════════════════════════════
% ML-04: Musterlösung Stromdetektive im Haushalt
% Physik Klasse 8 | Projektarbeit
% ═══════════════════════════════════════════════════════════════

\documentclass[11pt, a4paper]{article}

\usepackage[utf8]{inputenc}
\usepackage[T1]{fontenc}
\usepackage[ngerman]{babel}

\usepackage{helvet}
\renewcommand{\familydefault}{\sfdefault}

\usepackage{microtype}
\usepackage[margin=1.5cm, top=1.8cm, bottom=1.5cm]{geometry}
\usepackage{enumitem}
\usepackage{tabularx}
\usepackage{booktabs}
\usepackage{array}

\usepackage{amsmath, amssymb}
\usepackage{siunitx}
\sisetup{locale = DE, per-mode = symbol}

\usepackage{fancyhdr}
\usepackage{lastpage}

\usepackage{xcolor}
\usepackage{tcolorbox}
\tcbuselibrary{skins}

% Farben
\definecolor{physik}{HTML}{2c5aa0}
\definecolor{hellgrau}{HTML}{F5F5F5}
\definecolor{mittelgrau}{HTML}{888888}
\definecolor{loesung}{HTML}{c62828}

\colorlet{fachfarbe}{physik}

% Kopf- und Fußzeile
\pagestyle{fancy}
\fancyhf{}
\renewcommand{\headrulewidth}{0.4pt}
\lhead{\small Physik Klasse 8 | Stromdetektive -- \textcolor{loesung}{\textbf{MUSTERLÖSUNG}}}
\rhead{\small Lehrerexemplar}
\cfoot{\small\textcolor{mittelgrau}{Seite \thepage{} von \pageref{LastPage}}}

% Boxen
\newtcolorbox{loesungskasten}{
    colback=white,
    colframe=loesung,
    boxrule=1.5pt,
    arc=0pt,
    left=8pt, right=8pt, top=6pt, bottom=6pt,
    before skip=6pt, after skip=6pt
}

\setlength{\parindent}{0pt}
\setlength{\parskip}{4pt}

\begin{document}

% ═══════════════════════════════════════════════════════════════
% HEADER
% ═══════════════════════════════════════════════════════════════

\begin{center}
{\Large\bfseries\textcolor{loesung}{MUSTERLÖSUNG: Stromdetektive im Haushalt}}\\[0.3em]
{\normalsize Beispiellösung für einen typischen 3-Personen-Haushalt}
\end{center}
\vspace{-0.3em}
\hrule
\vspace{0.8em}

% ═══════════════════════════════════════════════════════════════
% BEISPIEL-TABELLE
% ═══════════════════════════════════════════════════════════════

\textbf{Geräte-Erfassung (Beispiel):}

\vspace{0.3em}
\small
\begin{center}
\renewcommand{\arraystretch}{1.3}
\begin{tabularx}{\textwidth}{|c|X|c|c|c|}
\hline
\textbf{Raum} & \textbf{Gerät} & \textbf{P (W)} & \textbf{t (h/Tag)} & \textbf{E (Wh/Tag)} \\
\hline
\hline
\multirow{3}{*}{Küche} & Kühlschrank & 15 & 24 & \textcolor{loesung}{360} \\
\cline{2-5}
 & Wasserkocher & 2000 & 0,1 & \textcolor{loesung}{200} \\
\cline{2-5}
 & Geschirrspüler & 1800 & 1 & \textcolor{loesung}{1800} \\
\hline
\multirow{3}{*}{Wohnzimmer} & Fernseher & 100 & 4 & \textcolor{loesung}{400} \\
\cline{2-5}
 & Spielkonsole & 120 & 1 & \textcolor{loesung}{120} \\
\cline{2-5}
 & Router & 10 & 24 & \textcolor{loesung}{240} \\
\hline
\multirow{2}{*}{Schlafzimmer} & LED-Lampe & 10 & 2 & \textcolor{loesung}{20} \\
\cline{2-5}
 & Handy-Ladegerät & 5 & 2 & \textcolor{loesung}{10} \\
\hline
\multirow{2}{*}{Bad} & Föhn & 2000 & 0,15 & \textcolor{loesung}{300} \\
\cline{2-5}
 & Elektr. Zahnbürste & 5 & 0,1 & \textcolor{loesung}{0,5} \\
\hline
\multirow{2}{*}{Arbeitszimmer} & Computer & 80 & 4 & \textcolor{loesung}{320} \\
\cline{2-5}
 & Monitor & 30 & 4 & \textcolor{loesung}{120} \\
\hline
\hline
\multicolumn{4}{|r|}{\textbf{Summe:}} & \textcolor{loesung}{\textbf{3890,5 Wh}} \\
\hline
\end{tabularx}
\end{center}

\normalsize
\vspace{0.5em}

% ═══════════════════════════════════════════════════════════════
% AUFGABE 1
% ═══════════════════════════════════════════════════════════════

\textbf{Aufgabe 1 -- Die größten Verbraucher (4 BE):}

\begin{loesungskasten}
\begin{tabular}{|c|l|c|l|}
\hline
\textbf{Platz} & \textbf{Gerät} & \textbf{E (Wh/Tag)} & \textbf{Warum so viel?} \\
\hline
1 & Geschirrspüler & 1800 & Hohe Leistung (1800 W), lange Laufzeit \\
\hline
2 & Fernseher & 400 & Mittlere Leistung, aber lange Nutzung (4 h) \\
\hline
3 & Kühlschrank & 360 & Niedrige Leistung, aber 24 h Betrieb \\
\hline
\end{tabular}

\vspace{0.3em}
\textit{Hinweis: Die Reihenfolge variiert je nach Haushalt. Typische Top-Verbraucher sind Kühl-/Gefriergeräte, Waschmaschine, Trockner, TV, PC.}
\end{loesungskasten}

% ═══════════════════════════════════════════════════════════════
% AUFGABE 2
% ═══════════════════════════════════════════════════════════════

\textbf{Aufgabe 2 -- Stromkosten berechnen (7 BE):}

\begin{loesungskasten}
\textbf{a) Tagesverbrauch:}

$E_{\text{Tag}} = 3890{,}5\,\text{Wh} = \textcolor{loesung}{\mathbf{3{,}89\,\text{kWh}}}$

\vspace{0.5em}
\textbf{b) Monatskosten:}

$\text{Kosten}_{\text{Monat}} = E_{\text{Tag}} \cdot 30 \cdot \text{Strompreis}$

$= 3{,}89\,\text{kWh/Tag} \cdot 30\,\text{Tage} \cdot 0{,}30\,\text{€/kWh}$

$= 116{,}7\,\text{kWh} \cdot 0{,}30\,\text{€/kWh} = \textcolor{loesung}{\mathbf{35{,}01\,€}}$

\vspace{0.5em}
\textbf{c) Jahreskosten:}

$\text{Kosten}_{\text{Jahr}} = E_{\text{Tag}} \cdot 365 \cdot \text{Strompreis}$

$= 3{,}89\,\text{kWh/Tag} \cdot 365\,\text{Tage} \cdot 0{,}30\,\text{€/kWh} = \textcolor{loesung}{\mathbf{426\,€}}$

\vspace{0.3em}
\textit{Hinweis: Ein typischer 3-Personen-Haushalt verbraucht ca. 3000--4000 kWh/Jahr, also 900--1200 €/Jahr. Der Wert hängt stark von Heizung/Warmwasser ab!}
\end{loesungskasten}

% ═══════════════════════════════════════════════════════════════
% AUFGABE 3
% ═══════════════════════════════════════════════════════════════

\textbf{Aufgabe 3 -- Sparmaßnahmen (5 BE):}

\begin{loesungskasten}
\begin{tabular}{|l|p{10cm}|}
\hline
\textbf{Gerät} & \textbf{Sparmaßnahme} \\
\hline
Geschirrspüler & Eco-Programm nutzen (spart bis zu 30\%), nur voll beladen starten \\
\hline
Fernseher & Kürzere Nutzung, Helligkeit reduzieren, Standby vermeiden \\
\hline
Kühlschrank & Temperatur auf 7°C einstellen, Tür schnell schließen, regelmäßig abtauen \\
\hline
\end{tabular}

\vspace{0.3em}
\textbf{Weitere akzeptable Maßnahmen:}
\begin{itemize}[leftmargin=1.5em, itemsep=2pt]
    \item Waschmaschine: 30°C statt 60°C, volle Trommel
    \item Beleuchtung: LED statt Glühbirne
    \item Standby: Steckerleiste mit Schalter
    \item Computer: Energiesparmodus, bei Nichtnutzung ausschalten
    \item Föhn: Haare an der Luft trocknen lassen
\end{itemize}
\end{loesungskasten}

% ═══════════════════════════════════════════════════════════════
% AUFGABE 4
% ═══════════════════════════════════════════════════════════════

\textbf{Aufgabe 4 -- Einsparpotenzial berechnen (4 BE):}

\begin{loesungskasten}
\textbf{Beispiel: Geschirrspüler mit Eco-Programm}

\vspace{0.3em}
\textbf{Vorher:} $P = 1800\,\text{W}$, $t = 1\,\text{h}$, $E = 1800\,\text{Wh/Tag}$

\textbf{Nachher:} $P = 1800\,\text{W}$, $t = 1\,\text{h}$, aber Eco spart 30\%: $E = 1260\,\text{Wh/Tag}$

\vspace{0.3em}
\textbf{Ersparnis pro Tag:}

$\Delta E = 1800 - 1260 = \textcolor{loesung}{\mathbf{540\,\text{Wh/Tag}}}$

\vspace{0.3em}
\textbf{Ersparnis pro Jahr:}

$\Delta E_{\text{Jahr}} = 540\,\text{Wh} \cdot 365 = 197.100\,\text{Wh} = \textcolor{loesung}{\mathbf{197{,}1\,\text{kWh}}}$

\vspace{0.3em}
\textbf{Geldersparnis:}

$197{,}1\,\text{kWh} \cdot 0{,}30\,\text{€/kWh} = \textcolor{loesung}{\mathbf{59{,}13\,€/Jahr}}$

\vspace{0.5em}
\textit{Hinweis: Die Berechnung ist vereinfacht. In der Realität variiert der Verbrauch, und das Eco-Programm läuft länger, verbraucht aber weniger Energie.}
\end{loesungskasten}

\vspace{1em}

% ═══════════════════════════════════════════════════════════════
% BEWERTUNGSHINWEISE
% ═══════════════════════════════════════════════════════════════

\begin{tcolorbox}[colback=hellgrau, colframe=physik, title=\textbf{Bewertungshinweise}, fonttitle=\bfseries\color{white}]

\textbf{Häufige Fehler:}
\begin{itemize}[leftmargin=1.5em, itemsep=2pt]
    \item Verwechslung von Wh und kWh (Faktor 1000)
    \item Nutzungsdauer in Minuten statt Stunden angegeben
    \item Unrealistische Leistungswerte (z.B. Kühlschrank mit 100 W statt 15 W Durchschnitt)
    \item Formel falsch angewandt: E = P · t, nicht E = P / t
\end{itemize}

\textbf{Akzeptanzbereiche:}
\begin{itemize}[leftmargin=1.5em, itemsep=2pt]
    \item Tagesverbrauch: 2--8 kWh realistisch
    \item Monatskosten: 20--80 € realistisch (ohne elektr. Heizung)
    \item Jahreskosten: 250--1000 € realistisch (ohne elektr. Heizung)
\end{itemize}

\textbf{Bewertungsprinzip:} Folgefehler berücksichtigen! Wenn die Tabelle falsch ist, aber die Folgeberechnungen mit den falschen Werten korrekt durchgeführt werden, gibt es trotzdem Punkte für die Rechenmethode.

\end{tcolorbox}

\end{document}
